\section{What Maintains Convergence and What Grows with $L$}
\label{sec:convergence-analysis-L}

This section provides a detailed analysis of what maintains convergence when stacking low-rank layers and what quantities grow with depth $L$.

\subsection{What Maintains Convergence}

\textbf{The convergence is maintained by three key boundedness conditions:}

\begin{enumerate}
\item \textbf{Bounded mixing matrices}: For each layer $i = 2, \ldots, L-1$:
\[
\|L^{(i)}\|_{\infty,1} = \sup_{c_i}\sum_{k=1}^r |L^{(i)}_{c_i,k}| \le rK.
\]
This ensures that each layer's forward propagation is bounded: $|H_i| \le rK \max_k |m_k^{(i)}|$.

\item \textbf{Bounded activations}: For each layer $i = 1, \ldots, L$:
\[
|\varphi_i(x)| \le K, \quad |\varphi_i'(x)| \le K, \quad |\varphi_i(x) - \varphi_i(y)| \le K|x-y|.
\]
This ensures that activation functions don't amplify errors unboundedly.

\item \textbf{Bounded data}: 
\[
|X| \le K \quad \text{almost surely}.
\]
This ensures that all expectations remain bounded.
\end{enumerate}

\textbf{The frozen first layer provides a stable base:} Since $f_1(C_1)$ is fixed and independent of $W(t)$, we have:
\[
H_1(t, c_1; X, W) = \langle f_1(c_1), X\rangle \quad \text{(independent of $W(t)$)}.
\]
This eliminates recursion at the base case and ensures errors don't compound from the input.

\subsection{What Grows with $L$: Error Propagation Analysis}

When stacking $L$ layers, errors propagate in two directions:

\paragraph*{Forward Propagation (Layer $i$ to $i+1$):}

For $i = 2, \ldots, L$, the pre-activation $H_i$ depends on the previous layer:
\[
H_i(t, c_i; X, W) = \sum_{k=1}^r L^{(i)}_{c_i,k}\, m_k^{(i)}(t; X, W),
\]
where for $i \ge 3$:
\[
m_k^{(i)}(t; X, W) = \mathbb{E}_{C_{i-1}}[w_{i-1}(t, C_{i-1}, k)\,\varphi_{i-1}(H_{i-1}(t, C_{i-1}; X, W))].
\]

\textbf{Error propagation bound:} If $|H_{i-1}(W') - H_{i-1}(W'')| \le \delta_{i-1}$, then:
\begin{align*}
|H_i(W') - H_i(W'')| &\le \sum_{k=1}^r |L^{(i)}_{c_i,k}|\,|m_k^{(i)}(W') - m_k^{(i)}(W'')|\\
&\le rK \max_k \mathbb{E}_{C_{i-1}}[|w_{i-1}'(C_{i-1}, k) - w_{i-1}''(C_{i-1}, k)|]\\
&\quad + rK^2 \max_k \mathbb{E}_{C_{i-1}}[|H_{i-1}'(C_{i-1}) - H_{i-1}''(C_{i-1})|]\\
&\le rK \mathscr{D}_t(W', W'') + rK^2 \delta_{i-1}.
\end{align*}

\textbf{Key observation:} The error in $H_i$ depends on:
\begin{itemize}
\item The weight difference $\mathscr{D}_t(W', W'')$ (bounded by our distance measure)
\item The error from the previous layer $\delta_{i-1}$ (multiplied by $rK^2$)
\end{itemize}

\paragraph*{Backward Propagation (Layer $i+1$ to $i$):}

For $i = 2, \ldots, L-1$, the backpropagated signal $B_k^{(i)}$ depends on the next layer:
\[
B_k^{(i)}(t; X, W) = \mathbb{E}_{C_i}[L^{(i)}_{C_i,k}\,\varphi_i'(H_i(t, C_i; X, W))\,w_i(t, C_i, k)].
\]

\textbf{Error propagation bound:} If $|B_k^{(i+1)}(W') - B_k^{(i+1)}(W'')| \le \delta_{B,i+1}$, then the update for layer $i$ depends on this error. However, the key is that $B_k^{(i)}$ itself depends on $H_i$ and $w_i$, which are bounded.

\subsection{The Gronwall Constant: Why $(1+rK)^{L-2}$ Works}

The critical insight is in the Gronwall argument. When we bound the update differences, we get:

\[
\left|\frac{\partial}{\partial t}\tilde{w}_i - \frac{\partial}{\partial t}w_i\right| \le K_t\left(\mathbb{E}_Z[|q_i|] + \mathbb{E}_Z[|q_{B,k}^{(i+1)}|]\right) + K_t\mathscr{D}_t(W, \tilde{W}).
\]

Each term contributes:
\begin{itemize}
\item $|q_i| \le rK \max_k(Q_{i,1,k} + Q_{i,2,k}) \le rK(\mathscr{D}_t + \text{concentration error})$
\item $|q_{B,k}^{(i+1)}| \le rK^2(\mathscr{D}_t + \text{concentration error})$ (from the next layer)
\end{itemize}

\textbf{The key is that these errors don't compound exponentially because:}
\begin{enumerate}
\item Each layer's contribution is \textbf{additive} in the Gronwall inequality, not multiplicative
\item The forward error $q_i$ is bounded by $rK \times \mathscr{D}_t$ (not $rK \times \delta_{i-1}$)
\item The backward error $q_{B,k}^{(i+1)}$ is bounded by $rK^2 \times \mathscr{D}_t$ (not $rK^2 \times \delta_{B,i+1}$)
\end{enumerate}

\textbf{Why this works:} We bound everything in terms of the global distance $\mathscr{D}_t(W, \tilde{W})$, not in terms of layer-wise errors that compound. This is the crucial difference from naive error propagation.

\subsection{The Accumulation Pattern}

When we apply Gronwall's lemma, we get:
\[
\mathscr{D}_t(W, \tilde{W}) \le K_T(1+rK)\int_0^t \mathscr{D}_s(W, \tilde{W}) ds + \text{(polynomial terms)}.
\]

For $L$ layers, each intermediate layer contributes a factor, but \textbf{the key is that they contribute to the same Gronwall constant}, not separate ones. The accumulation is:

\[
\text{Gronwall constant} = K_T \prod_{i=2}^{L-1} (1 + \|L^{(i)}\|_{\infty,1}) \le K_T(1+rK)^{L-2}.
\]

\textbf{Why this is polynomial, not exponential:}
\begin{itemize}
\item Each factor $(1+rK)$ is \textbf{bounded} (doesn't grow with $L$)
\item The number of factors is $L-2$ (linear in $L$)
\item The product is $(1+rK)^{L-2}$, which is \textbf{polynomial in $L$} (specifically, exponential in $\log(1+rK) \times L$, but since $rK$ is fixed, this is polynomial)
\end{itemize}

\subsection{What Would Break Convergence}

Convergence would fail if:

\begin{enumerate}
\item \textbf{Unbounded mixing matrices}: If $\|L^{(i)}\|_{\infty,1}$ grew with $L$ or $i$, then the Gronwall constant would grow exponentially.

\item \textbf{Unbounded activations}: If activations could amplify errors unboundedly (e.g., exponential activations), then errors would compound exponentially.

\item \textbf{Error compounding}: If we bounded errors in terms of previous layer errors rather than the global distance $\mathscr{D}_t$, then errors would compound as $\delta_i \le (rK)^i \delta_1$, which grows exponentially.

\item \textbf{No frozen base}: If the first layer were trainable and could diverge, then all subsequent layers would inherit this instability.
\end{enumerate}

\subsection{The Key Insight: Global Distance Bounding}

\textbf{The crucial technique that maintains convergence is bounding everything in terms of the global distance $\mathscr{D}_t(W, \tilde{W})$ rather than layer-wise errors.}

Instead of:
\[
\delta_i \le (rK)^i \delta_1 \quad \text{(exponential growth)}
\]

We use:
\[
\delta_i \le rK \mathscr{D}_t \quad \text{(bounded by global distance)}
\]

This ensures that all layer errors are controlled by the same global quantity, which we bound using Gronwall's lemma with a constant that grows only polynomially in $L$.

\subsection{Conclusion}

\textbf{What maintains convergence:}
\begin{itemize}
\item Bounded mixing matrices: $\|L^{(i)}\|_{\infty,1} \le rK$ (fixed, independent of $L$)
\item Bounded activations: $|\varphi_i|, |\varphi_i'| \le K$ (fixed)
\item Bounded data: $|X| \le K$ (fixed)
\item Frozen first layer: Provides stable base
\item Global distance bounding: All errors bounded in terms of $\mathscr{D}_t$, not layer-wise
\end{itemize}

\textbf{What grows with $L$:}
\begin{itemize}
\item Gronwall constant: $(1+rK)^{L-2}$ (polynomial in $L$, not exponential)
\item Number of layers to bound: $L-2$ (linear in $L$)
\item Logarithmic factors: $\log(n_{\max}^2 r^{L-2}/\delta)$ (polynomial in $L$)
\end{itemize}

\textbf{The key is that growth is polynomial, not exponential, because:}
\begin{enumerate}
\item Each layer contributes a \textbf{bounded factor} $(1+rK)$
\item These factors \textbf{multiply} (not compound exponentially)
\item All errors are bounded in terms of the \textbf{global distance} $\mathscr{D}_t$, not layer-wise errors
\end{enumerate}

This is why the proof works for arbitrary depth $L$: the boundedness conditions ensure polynomial growth, and the global distance bounding prevents exponential error compounding.
